% Laget automatisk av elevark/lag_elevark.py fra sim/skratt-kast.html. Endringer her blir overskrevet.
\documentclass[11pt]{article}
\usepackage{elevark}
\begin{document}
\elevarkhode{Fysikk 2 · Bevegelse og krefter i to dimensjoner}{Skrått kast}{https://lektorbo.no/sim/skratt-kast.html}{Simuleringen}

En ball som kastes skrått, beveger seg i to retninger samtidig. Vannrett er farten konstant, og loddrett virker tyngdeakselerasjonen. Hvilken vinkel gir lengst kast? Og hva gjør luftmotstanden med banen?

\section*{Oppgaver}

Start med $v_{0} = 15\,\mathrm{m/s}$, $α = 45^{\circ}$ og $h_{0} = 0$, uten luftmotstand.

\begin{oppgaveliste}
\oppgave \fase{gjett}{Gjett}Hvordan endrer $v_{x}$ og $v_{y}$ seg under kastet? Kast og se på pilene og tallene.
\svar
\oppgave \fase{test}{Regn}Dekomponer utgangsfarten. Regn ut tiden ballen er i lufta, rekkevidden og den største høyden. Sjekk.
\svar
\oppgave \fase{gjett}{Gjett}Kast med 30° og så med 60°. Hvilket kast går lengst? Hvilket går høyest og varer lengst? Forklar.
\svar
\oppgave \fase{test}{Undersøk}Finn vinkelen som gir lengst kast når $h_{0} = 0$. Vis at rekkevidden er $R = v_{0}^{2} \sin (2α) / g$.
\svar
\oppgave \fase{test}{Undersøk}Sett $h_{0} = 10\,\mathrm{m}$. Er 45° fortsatt den beste vinkelen? Prøv 35°, 40° og 45°. Hvor stor er farten når ballen lander, og avhenger den av vinkelen?
\svar
\oppgave \fase{forklar}{Forklar}Slå på luftmotstand og kast med 25~m/s og 45°. Sammenlign med kastet uten luftmotstand. Hvorfor er banen ikke lenger symmetrisk?
\svar
\oppgave \fase{utfordring}{Utfordring}En ball kastes vannrett fra 10~m høyde. Hvor fort må den kastes for å lande 15~m ut? Test med $α = 0^{\circ}$.
\svar
\end{oppgaveliste}

\end{document}
